\begin{table}[!h]
\centering
\renewcommand{\arraystretch}{0.98}
\setlength{\tabcolsep}{8pt}
\captionsetup{justification=centering}
\caption{Heterogeneous Effects on Migration Intention}
\begin{tabular*}{\textwidth}{l@{\extracolsep{\fill}}cccc}
\hline
& \multicolumn{2}{c}{Pop.\ density} & \multicolumn{2}{c}{Female pop.} \\
\cmidrule(l){2-3} \cmidrule(l){4-5}
& Below & Above & Below & Above \\
\hline
Spanish share 1936-1955$\times$Post & 0.037** & -0.004 & 0.032* & -0.002 \\
 & (0.018) & (0.041) & (0.018) & (0.042) \\
Spanish share 1956-1978$\times$Post & -0.009 & 0.065 & 0.004 & 0.064 \\
  & (0.088) & (0.126) & (0.094) & (0.131) \\
\addlinespace
Observations & 1,240 & 6,208 & 1,997 & 5,451 \\
$p$-value ($\beta_{36{-}55} = \beta_{56{-}78}$) & 0.654 & 0.676 & 0.795 & 0.697 \\
\hline
\end{tabular*}
\begin{tabular*}{\textwidth}{l@{\extracolsep{\fill}}cccc}
& \multicolumn{2}{c}{Mean years educ.} & \multicolumn{2}{c}{Median age} \\
\cmidrule(l){2-3} \cmidrule(l){4-5}
& Below & Above & Below & Above \\
\hline
Spanish share 1936-1955$\times$Post & -0.016 & -0.028 & -0.002 & -0.049 \\
 & (0.040) & (0.044) & (0.039) & (0.045) \\
Spanish share 1956-1978$\times$Post & 0.403* & 0.102 & 0.403* & 0.123 \\
  & (0.237) & (0.129) & (0.241) & (0.130) \\
\addlinespace
Observations & 1,955 & 5,493 & 3,095 & 4,353 \\
$p$-value ($\beta_{36{-}55} = \beta_{56{-}78}$) & 0.128 & 0.445 & 0.145 & 0.316 \\
\hline
\end{tabular*}
\begin{tabular*}{\textwidth}{l@{\extracolsep{\fill}}cccc}
& \multicolumn{2}{c}{Share in labor force} & \multicolumn{2}{c}{Share unemployed} \\
\cmidrule(l){2-3} \cmidrule(l){4-5}
& Below & Above & Below & Above \\
\hline
Spanish share 1936-1955$\times$Post & 0.028 & -0.014 & 0.046 & 0.005 \\
 & (0.039) & (0.041) & (0.047) & (0.033) \\
Spanish share 1956-1978$\times$Post & 0.168 & 0.042 & -0.135 & 0.108 \\
  & (0.219) & (0.121) & (0.177) & (0.100) \\
\addlinespace
Observations & 1,705 & 5,743 & 2,112 & 5,336 \\
$p$-value ($\beta_{36{-}55} = \beta_{56{-}78}$) & 0.588 & 0.727 & 0.412 & 0.420 \\
\hline
\end{tabular*}
\begin{tabular*}{\textwidth}{l@{\extracolsep{\fill}}cccc}
& \multicolumn{2}{c}{Left vote share} & \multicolumn{2}{c}{Ideological alternation} \\
\cmidrule(l){2-3} \cmidrule(l){4-5}
& Below & Above & Below & Above \\
\hline
Spanish share 1936-1955$\times$Post & 0.028 & 0.022 & 0.018 & 0.009 \\
 & (0.049) & (0.034) & (0.030) & (0.057) \\
Spanish share 1956-1978$\times$Post & -0.074 & 0.073 & 0.106 & -0.014 \\
  & (0.174) & (0.103) & (0.094) & (0.192) \\
\addlinespace
Observations & 3,750 & 3,698 & 4,082 & 3,366 \\
$p$-value ($\beta_{36{-}55} = \beta_{56{-}78}$) & 0.644 & 0.701 & 0.454 & 0.925 \\
\hline
\end{tabular*}
\addvspace{0.2em}
\captionsetup{font=footnotesize, justification=justified, singlelinecheck=false}
\caption*{\footnotesize \textit{Notes}: The dependent variable is migration intention. Each column reports estimates for the subsample of municipalities below or above the median of the specified municipal-level variable. Medians are defined using the full sample of 312 municipalities, and the corresponding median cutoffs used to split the sample are also calculated from this full municipal sample. The split variables are defined as follows. Population density, the female share of the population, mean years of education, median age, the labor-force participation rate among individuals aged 14 and older, and the unemployment rate among individuals aged 14 and older are measured using the 2010 census. Broad-left vote share is the average share of valid votes received by left and center-left parties over the pre-reform elections held between 2011 and 2021. Ideological alternation is an indicator equal to one if the ideological bloc of the winning party changes relative to the previous election of the same type; the split variable is the municipality-level average of this indicator over the 2011--2021 pre-reform period. The median value used to split the sample for each variable is: Pop.\ density (13.000); Pop.\ density (14596.960); Female pop. (50.540); Female pop. (53.968); Mean years educ. (7.583); Mean years educ. (10.923); Median age (28.000); Median age (38.000); Share in labor force (63.769); Share in labor force (79.671); Share unemployed (3.206); Share unemployed (5.991); Left vote share (63.722); Left vote share (88.801); Ideological alternation (40.000); Ideological alternation (100.000). All specifications include municipality and year fixed effects and individual controls for age and a male indicator. The reported coefficients are average marginal effects from a logit specification, expressed as the change in the probability of reporting migration intention per one-unit increase in the corresponding regressor, where one unit corresponds to one percentage point of the municipal Spanish-born population share. Standard errors are clustered at the municipality level and reported in parentheses (56 clusters). The row labeled $p$-value: $\beta^{1936-1955} = \beta^{1956-1978}$ reports the p-value from a two-sided Wald test of equality between the two cohort-specific coefficients within each subsample. Tests of equality of the coefficient on Spanish share $\times$ Post between the below-median and above-median subsamples ($H_0: \beta_{\text{Below}} = \beta_{\text{Above}}$), computed as z-statistics on the difference between the two coefficients estimated on independent subsamples, yield the following p-values, reported as ($\beta^{1936-1955}$ / $\beta^{1956-1978}$) for each variable: Pop.\ density (0.355 / 0.634); Female pop. (0.460 / 0.706); Mean years educ. (0.838 / 0.265); Median age (0.422 / 0.306); Share in labor force (0.457 / 0.616); Share unemployed (0.468 / 0.233); Left vote share (0.929 / 0.467); Ideological alternation (0.886 / 0.576). * $p<0.10$, ** $p<0.05$, *** $p<0.01$. \\textit{Sources}:  LAPOP AmericasBarometer, 2023 round (dependent variables and individual controls) National Electoral Commission (Left vote share, Ideological alternation); IPUMS Argentine Census 1970 and 1980 (Spanish immigration shares); IPUMS Argentine Census 2010 (municipal demographic and socioeconomic characteristics).}
\end{table}
